\section{Conclusion and Outlook}
\label{sec:conclusion_outlook}
\ihead{Conclusion and Outlook}
This project developed a \ac{ML} pipeline for short-term electricity demand forecasting in the 50Hertz control area. Starting from a Seasonal Na\"{i}ve baseline with a MAPE of 8.57\%, the final model, a tuned and refitted \acs{XGB} estimator embedded in a recursive forecasting framework, achieved a MAPE of 2.94\% on the 2024 test year. The inclusion of exogenous features proved to be the single most impactful factor, roughly halving the forecasting error compared to models relying solely on lagged demand values. Feature importance and \ac{SHAP} analysis confirmed that the model relies primarily on recent demand observations, cyclical temporal encodings, and year-over-year load patterns, while solar generation forecasts contribute a physically interpretable inverse effect. The temporal error analysis revealed that forecasting accuracy is highest during nighttime and midweek periods and degrades during holiday periods, weekends and later hours of the forecast horizon. Several limitations should be noted. The weather data used in this study consists of realized observations rather than actual meteorological forecasts, meaning that the reported accuracy represents an optimistic estimate of operational performance. Furthermore, the backtesting framework generates forecasts exclusively from midnight, which prevents a clean separation between hour-of-day effects and recursive error accumulation. The analysis is also restricted to a single control area, and the models produce only deterministic point forecasts without quantifying prediction uncertainty. Future work could address these limitations in several directions. Incorporating actual weather forecasts instead of realized observations would provide a more realistic assessment of operational accuracy. Extending the analysis to additional control areas would test the generalizability of the proposed approach. Probabilistic forecasting methods, such as conformal prediction intervals~\cite{zaffran_2022_adaptive} or quantile regression~\cite{hong_2016_probabilistic}, could complement the point forecasts with uncertainty estimates that are valuable for risk-aware grid operation. Finally, a comparison with deep learning architectures such as Temporal Fusion Transformers could reveal whether the remaining forecasting error can be further reduced by capturing more complex temporal dependencies~\cite{lim_2021_temporal}.